% 3 x 2 meshes on [0,3] x [0,2] of the patch test, the eigenvalue study and the
% transient study (studies/ms_spectrum.m): (a) curved 0.15, (b) curved 0.20,
% (c) angular 0.3 + curved 0.15 (usemDistortMesh 'combined', 0.30).
% Thick boundary line (hatched in (a)): u prescribed at all boundary nodes.
% data: mesh32-<mesh>.dat (studies/ms_mesh_figures.m)
\tracinglostchars=0 % blank lines (jumps) in the mesh data
\begin{tikzpicture}[plabel/.style={anchor=north, font=\small, inner sep=0pt}]
\begin{groupplot}[group style={group name=m32, group size=3 by 1, horizontal sep=0.9cm},
    hide axis, clip=false, x=1.62cm, y=1.62cm, xmin=0, xmax=3, ymin=0, ymax=2,
    mesh/.style={black, line width=0.4pt},
    bnd/.style={black, line width=1.2pt}]
  \nextgroupplot
    \addplot[mesh] table {\data/mesh32-curved15.dat};
    \addplot[bnd] coordinates {(0,0) (3,0) (3,2) (0,2) (0,0)};
    % hatching (clamped boundary)
    \pgfplotsinvokeforeach{0.1,0.25,...,3.0}{
      \draw[line width=0.3pt] (axis cs:#1,0) -- (axis cs:#1-0.1,-0.1);
      \draw[line width=0.3pt] (axis cs:#1-0.1,2) -- (axis cs:#1,2.1);}
    \pgfplotsinvokeforeach{0.1,0.25,...,2.0}{
      \draw[line width=0.3pt] (axis cs:0,#1-0.1) -- (axis cs:-0.1,#1);
      \draw[line width=0.3pt] (axis cs:3,#1) -- (axis cs:3.1,#1-0.1);}
    \node[plabel] at (axis cs:1.5,-0.22) {(a) curved, $0.15$};
  \nextgroupplot
    \addplot[mesh] table {\data/mesh32-curved20.dat};
    \addplot[bnd] coordinates {(0,0) (3,0) (3,2) (0,2) (0,0)};
    \node[plabel] at (axis cs:1.5,-0.22) {(b) curved, $0.20$};
  \nextgroupplot
    \addplot[mesh] table {\data/mesh32-combined30.dat};
    \addplot[bnd] coordinates {(0,0) (3,0) (3,2) (0,2) (0,0)};
    \node[plabel] at (axis cs:1.5,-0.22) {(c) angular $0.3$ + curved $0.15$};
\end{groupplot}
\node[anchor=north, font=\footnotesize, yshift=-0.75cm] at (m32 c2r1.south)
  {thick (hatched) boundary: $\vect u$ prescribed at all boundary nodes (patch test, smooth field);
   $\vect u=\vect 0$ (eigenvalue and transient study)};
\end{tikzpicture}
